\documentclass[11pt]{article}
\usepackage[utf8]{inputenc}
\usepackage[T1]{fontenc}
\usepackage{lmodern}
\usepackage{amsmath,amssymb}
\usepackage{graphicx}
\usepackage{booktabs}
\usepackage{ulem}
\usepackage{fixltx2e}
\usepackage[margin=20mm,a4paper]{geometry}
\usepackage{hyperref}
\hypersetup{colorlinks=true,linkcolor=blue,urlcolor=blue}
\usepackage{kotex}
\title{MDCV(1)}

\date{September 22, 2026}

\begin{document}
\maketitle

\section*{NAME}\label{name}

\textbf{mdcv} --- My Document Converter 의 mdoc 샘플

\section*{SYNOPSIS}\label{synopsis}

\textbf{} [Fl f Ar format] [Fl t Ar format] \emph{file}

\section*{DESCRIPTION}\label{description}

mdoc 은 BSD 매뉴얼 페이지 형식입니다. \emph{기울임} 과 \textbf{굵게} 와 \verb|코드| 를 씁니다.
This sentence is in English.

옵션은 다음과 같습니다.

\begin{description}
\item[\textbf{-f Ar format}] 입력 형식을 지정합니다.
\item[\textbf{-t Ar format}] 출력 형식을 지정합니다.
\end{description}

\section*{EXAMPLES}\label{examples}

\begin{itemize}
\item 첫 번째 항목
\item 두 번째 항목
\item 세 번째 항목
\end{itemize}

\begin{verbatim}
mdcv -f markdown -t html sample.md
\end{verbatim}

\section*{SEE ALSO}\label{see-also}

\textbf{pandoc(1)}

\end{document}
